\documentclass[aspectratio=169,11pt]{beamer}
\usetheme{default}
\usecolortheme{default}
\setbeamertemplate{navigation symbols}{}
\setbeamertemplate{footline}[frame number]
\usepackage[T1]{fontenc}
\usepackage{booktabs}

\title{Agentic Workloads for SeBS}
\subtitle{Progress update}
\author{Elif \"Ozsoy \and Ahmet Efe Serdar \and Jan Staszewicz \and Ivan Bozhanin}
\institute{DPHPC HS26 \quad Mentor: Marcin Copik}
\date{19 November 2026}

\begin{document}

\begin{frame}
  \titlepage
\end{frame}

% ---------------------------------------------------------------------------
\begin{frame}{Goal}
  \begin{itemize}
    \item AI agents do not only generate text: they \textbf{call tools} (browse, run code, fetch pages).
    \item Tool calls are short, bursty and independent: a natural fit for \textbf{serverless functions}.
    \item Question: \emph{how do agent workloads behave on serverless platforms?}
    \item SeBS has 15 benchmarks, \textbf{none of them agentic}.
  \end{itemize}
  \vspace{1em}
  Deliverable: a collection of agentic benchmarks for SeBS (single functions and pipelines),
  reproducible, runnable on any platform, without live LLM inference.
\end{frame}

% ---------------------------------------------------------------------------
\begin{frame}[fragile]{Method: record and replay}
\footnotesize
\begin{verbatim}
RECORD (once)    agent loop -> LLM decides -> tool runs for real
                 log every LLM turn and tool call       -> trace.json

REPLAY (always)  trace.json -> wait LLM time / call real tool
                 no LLM, identical every run            -> results.json
\end{verbatim}
\normalsize
  \begin{itemize}
    \item One trace format and one replayer shared by all benchmarks.
    \item Inputs are stored snapshots (pages, repositories, sites), never the live web.
    \item Replaying the same trace: \textbf{2.00\,s $\pm$ 0.03\,s} with snapshots;
          one run against live Wikipedia took 25\,s.
  \end{itemize}
\end{frame}

% ---------------------------------------------------------------------------
\begin{frame}[fragile]{Trace format (v0.1)}
\begin{verbatim}
{"task": "Compare serverless and cloud computing",
 "steps": [
   {"type": "llm",  "latency_s": 1.2, "input_tokens": 850},
   {"type": "tool", "tool": "web-fetch", "group": 1, "args": {"page": "..."}},
   {"type": "tool", "tool": "browser",   "group": 1, "args": {...}},
   {"type": "tool", "tool": "code-exec", "session": "ws", "expect": {...}},
   {"type": "tool", "tool": "pipeline",  "inputs_from": [5]}, ...]}
\end{verbatim}
  \begin{itemize}
    \item \texttt{group}: parallel tool calls in one LLM turn
    \item \texttt{session}: stateful tools share a workspace
    \item \texttt{inputs\_from}: output of one function feeds the next (via storage)
    \item \texttt{expect}: detects divergence from the recording
  \end{itemize}
\end{frame}

% ---------------------------------------------------------------------------
\begin{frame}{Four benchmarks (\texttt{700.agentic})}
  \small
  \begin{tabular}{@{}llp{7.6cm}@{}}
    \toprule
    Benchmark & Owner & What the agent does \\
    \midrule
    \texttt{710.web-fetch} & Ivan  & MCP fetch server: page $\to$ main content $\to$ Markdown $\to$ truncate \\
    \texttt{720.code-exec} & Ahmet & pinned repository: install, run tests, grep, patch, lint \\
    \texttt{730.browser}   & Jan   & headless Chromium on stored sites: form, Wikipedia, 2\,000-item shop \\
    \texttt{740.pipeline}  & Elif  & multi-step task, results handed off between functions via storage \\
    \bottomrule
  \end{tabular}
  \vspace{1em}
  \begin{itemize}
    \item Each: test/small/large inputs, output validation, own timing breakdown.
    \item All four replayed by the same runner in one agent task (see demo).
  \end{itemize}
\end{frame}

% ---------------------------------------------------------------------------
\begin{frame}{First results per benchmark (local, SeBS local backend)}
  \begin{itemize}
    \item \textbf{web-fetch:} CPU-bound, not I/O-bound. 3\,MB page: 2.4\,s parsing vs.\ 10\,ms fetch;
          the model receives 5\,000 of 260\,000 produced characters.
    \item \textbf{browser:} $\approx$475\,MB package, about $2\times$ the AWS Lambda zip limit (250\,MB);
          8 processes, 350--450\,MB memory; in a warm run half the time is closing Chromium.
    \item \textbf{code-exec:} a bug-fix session on \texttt{click} (7 tool calls) fails, then passes,
          exactly as recorded. Its state must survive between calls (next slide).
    \item \textbf{pipeline:} for a 1\,MB result, hand-off plus invocation overhead exceeds the actual work.
  \end{itemize}
\end{frame}

% ---------------------------------------------------------------------------
\begin{frame}{Stateful tools: keep the function alive, or save and load?}
  Bug-fix session on \texttt{click} (7 tool calls, 3 LLM turns between them):
  \vspace{0.5em}
  \begin{center}
  \begin{tabular}{@{}lrrr@{}}
    \toprule
    Mode & Invocations & Billed time & End-to-end \\
    \midrule
    whole session in one call (waits for the LLM) & 1 & 5.06\,s & 6.22\,s \\
    one call per step, workspace via storage       & 7 & 3.91\,s & 6.67\,s \\
    \bottomrule
  \end{tabular}
  \end{center}
  \begin{itemize}
    \item One call: \textbf{about a third of the billed time is idle waiting} for the model.
    \item Per step: \textbf{23\% cheaper, 7\% slower}, since every step saves and loads the workspace.
    \item Ahmet's measurement on fresh containers for every step: 11.7\,s,
          so per-step only pays off while containers stay warm.
  \end{itemize}
  \footnotesize Medians of 3 runs, local; per-step mode by Ahmet, replayed by the shared runner.
\end{frame}

% ---------------------------------------------------------------------------
\begin{frame}{Real traces with a local model}
  \begin{itemize}
    \item Recorder runs a tool-calling agent against any OpenAI-compatible API.
    \item Without API credits so far: \texttt{qwen2.5:3b} via Ollama on a laptop CPU.
    \item 28 traces (5 tasks $\times$ 5 runs, one failed, + 4 earlier):
      \begin{itemize}
        \item same task, different behaviour: one task always fetched 2 pages \textbf{in parallel},
              others mostly 1 page;
        \item page re-reads with \texttt{start\_index}, tool errors returned to the model.
      \end{itemize}
    \item Caveat: CPU inference makes LLM time $>96\%$ of every trace.
          Realistic latencies need a hosted model (Muse credits) or latency scaling.
  \end{itemize}
\end{frame}

% ---------------------------------------------------------------------------
\begin{frame}[fragile]{Demo: one agent task across all four benchmarks}
\begin{verbatim}
think 1.1s -> web-fetch || browser (parallel)            0.77 s
think 0.9s -> code-exec SESSION, 7 tool calls            3.76 s
                (2.20 s work + 1.55 s LLM wait inside the function)
think 0.7s -> pipeline produce -> consume (inputs_from)  0.06 s
think 1.4s -> answer
total 8.7 s over 3 runs, tools 53 %
\end{verbatim}
  \begin{itemize}
    \item LLM latencies here are synthetic ($\sim$1\,s per turn, typical for hosted APIs).
  \end{itemize}
\end{frame}

% ---------------------------------------------------------------------------
\begin{frame}{What agents do to serverless}
  \begin{itemize}
    \item \textbf{Agents are stateful, functions are not.} Workspace or browser tab is lost
          between calls: keep one invocation alive (idle billing) or save/load per step (cheaper, slower).
    \item \textbf{Parallel tool calls.} Locally they queue on one container;
          on a cloud platform they start new containers (cold starts).
    \item \textbf{Heavy tools.} Browser exceeds package limits; start-up dominates short actions.
    \item \textbf{Data hand-off} between tool functions can cost more than the work.
    \item Side finding: the SeBS local server's \texttt{begin}/\texttt{end} only time the
          function import, not the handler.
  \end{itemize}
\end{frame}

% ---------------------------------------------------------------------------
\begin{frame}{Next steps and questions}
  \textbf{Next}
  \begin{itemize}
    \item Cloud runs on AWS/GCP free tier; replay through SeBS triggers.
    \item Session modes on a real cloud: cold starts per step, storage latency.
    \item More traces; justify tool selection with AgentSysBench and tool-usage studies.
    \item One PR per benchmark.
  \end{itemize}
  \textbf{Questions}
  \begin{itemize}
    \item Muse credits for realistic traces?
    \item Should the replayer become a SeBS experiment?
    \item Snapshots: into the benchmarks-data repository?
  \end{itemize}
\end{frame}

\end{document}
